\chapter{Analysing the product of multiplying two quaternions \(PQ\).}

Before we proceed any further, let us first rewrite the final result from the previous section. However this time, as
part of the process of rewriting the final result, let us also highlight those terms which have certain particular
features in common with other terms. As you may have already noticed from the final result in the previous section --
or should hopefully notice in a moment, some of the terms share features with multiple other terms, while others share
features with only one other term.\\

Rewriting and highlighting the final result yields the following;

\[
    \begin{tikzpicture}
        \matrix (m)
        [
            matrix of math nodes,
            row sep=0pt,
            column sep=0pt
        ]
        {
            PQ = & |[name=m-1-1]| \, (P_0Q_0 \!\! & - & |[name=m-1-2]|P_1Q_1 & - & |[name=m-1-3]|P_2Q_2 & - & |[name=m-1-4]|P_3Q_3) & + &  \\
                 & |[name=m-2-1]|\quaternion{i}(P_0Q_1         & + & |[name=m-2-2]|P_1Q_0 & + & |[name=m-2-3]|P_2Q_3 & - & |[name=m-2-4]|P_3Q_2) & + &  \\
                 & |[name=m-3-1]|\quaternion{j}(P_0Q_2         & - & |[name=m-3-2]|P_1Q_3 & + & |[name=m-3-3]|P_2Q_0 & + & |[name=m-3-4]|P_3Q_1) & + & \\
                 & |[name=m-4-1]|\quaternion{k}(P_0Q_3         & + & |[name=m-4-2]|P_1Q_2 & - & |[name=m-4-3]|P_2Q_1 & + & |[name=m-4-4]|P_3Q_0)  \\
        };

        \draw[dashed, red, thin]
            (m-2-1.north west) rectangle (m-4-4.south east);

        \begin{scope}[on background layer]
        \draw[blue, line width=20pt, opacity=1.00, line cap=round]
           ($(m-2-2.center)+(0,0)$) -- ($(m-2-2.center)!1.05!(m-4-4.center)$);
        \end{scope}

        \begin{scope}[on background layer]Straight
        \draw[orange, line width=20pt, opacity=1.00, line cap=round]
            ($(m-2-3.center)+(0,0)$) -- ($(m-2-4.center)$);
        \end{scope}

        \begin{scope}[on background layer]
        \draw[green, line width=20pt, opacity=1.00, line cap=round]
            ($(m-3-2.center)+(0,0)$) -- ($(m-3-2.center)$);
        \end{scope}

        \begin{scope}[on background layer]
        \draw[purple, line width=20pt, opacity=1.00, line cap=round]
            ($(m-4-2.center)+(0,0)$) -- ($(m-4-3.center)$);
        \end{scope}

        \begin{scope}[on background layer]
        \draw[green, line width=20pt, opacity=1.00, line cap=round]
            ($(m-3-4.center)+(0,0)$) -- ($(m-3-4.center)$);
        \end{scope}

                \begin{scope}[on background layer]
        \draw[yellow, line width=20pt, opacity=1.00, line cap=round]
            ($(m-1-1.center)+(0,0)$) -- ($(m-1-4.center)$);
        \end{scope}

        \begin{scope}[on background layer]
        \draw[red, line width=20pt, opacity=1.00, line cap=round]
            ($(m-2-1.center)+(0,0)$) -- ($(m-4-1.center)$);
        \end{scope}

    \end{tikzpicture}
\]

Note that the final result is composed of sixteen different terms and that;

\begin{itemize}
    \item four of these terms make use of the quaternion unit \(1\), and
    \item the other twelve of these terms make use of one of the three imaginary quaternion units, i.e. \(\qi\),
          \(\qj\), or
          \(\quaternion{k}\).
\end{itemize}

The four terms which make use of the quaternion unit \(1\) are highlighted in yellow and when taken together, they act
as a scalar component within the final result.\\

The other twelve terms which make up the final result are surrounded by a thin dashed red line, and when taken together,
they act as a vector component within the final result.


\section{The scalar component of the quaternion product.}

As we said only a moment ago, the four terms which make use of the quaternion unit 1 are highlighted in yellow and when
taken together, they act as a scalar component within the final result. Upon inspection of these four terms, you should
be able to see that this scalar component can be defined as follows;

\vspace{0.25\baselineskip}

\begin{equation*}
    \text{Scalar component}(PQ) = P_{0}Q_{0} - \vec{P} \cdot \vec{Q}
\end{equation*}
\vspace{0.0\baselineskip}

where \(\vec{P} \cdot \vec{Q}\) denotes the dot, scalar, or inner product between the two vectors \(\vec{P}\) and
\(\vec{Q}\).


\section{The vector component of the quaternion product.}

Once again, as we said only a moment ago, the other twelve terms which make up the final result are surrounded by a thin
dashed red line, and when taken together, they act as a vector component within the final result. \\

\textbullet \; The term : \(P_{0}\vec{Q}\) \\

Within these twelve terms, we can see straight away that the three terms which have been highlighted in red can be
defined as;

\begin{equation*}
    P_{0}\vec{Q}
\end{equation*}
\vspace{0.0\baselineskip}

where;

\begin{equation*}
    \vec{Q} = \qi Q_1 + \qj Q_2 + \quaternion{k}Q_3
\end{equation*}
\vspace{0.0\baselineskip}

\textbullet \; The term : \(\vec{P}Q_{0}\) \\

Similarly, we can see that the three terms which have been highlighted in blue can be defined as;

\begin{equation*}
    \vec{P}Q_{0}
\end{equation*}

where;

\begin{equation*}
    \vec{P} = \qi P_1 + \quaternion{j} P_2 + \quaternion{k}P_3
\end{equation*}
\vspace{0.0\baselineskip}

\textbullet \; The term : \(\vec{P} \times \vec{Q}\) \\

Now all that's left, are the terms which are highlighted in orange, green, and purple. If you are familiar with
determinants, you might notice that we can generate these exact same terms using determinants and an appropriately
populated \(3 \times 3\) matrix. In order to do this, we would need a \(3 \times 3\) matrix which is populate as
follows;

\begin{equation*}
    \begin{vmatrix}
        \qi & \qj & \quaternion{k} \\
        P_1      & P_2      & P_3      \\
        Q_1      & Q_2      & Q_3
    \end{vmatrix}
\end{equation*}
\vspace{0.0\baselineskip}

Evaluating this determinant gives us what is known as the cross product of \(\vec{P}\) and \(\vec{Q}\), and which is
denoted mathematically as;

\begin{equation*}
    \vec{P} \times \vec{Q}
\end{equation*}
\vspace{0.0\baselineskip}

As a point of interest, note that the cross product doesn't contain any real terms. That is, it doesn't contain either
\(P_0\) or \(Q_0\). Keep this in mind, as we will return to this fact again later. \\

We won't discuss the cross product any further here, however we will return to it soon. \\

\textbullet \; Bringing these three terms together. \\

We have now accounted for all twelve of the terms which comprise the vector component of the final result. If we sum
together all three of the terms which we just derived, then we end up with following definition for our vector component.

\begin{equation*}
    \text{Vector component}(PQ) = P_0\vec{Q} + Q_0\vec{P} + (\vec{P}\times\vec{Q})
\end{equation*}

where;

\begin{itemize}
    \item \(P_0\vec{Q}\) corresponds to all of the elements in the final result which are highlighted in red.
    \item \(Q_0\vec{P}\) corresponds to all of the elements in the final result which are highlighted in blue.
    \item \(\vec{P}\times\vec{Q}\) corresponds to all of the elements in the final result which are highlighted in
                              orange, green, and purple.
\end{itemize}

From this, we can see that the vector component of the final result is comprised of two scalar products and one cross
product.

\section{Adding together the scalar and vector components of the quaternion product.}

If we bring together our definitions for the scalar part and the vector part of the final result, then we will get
the following result for our quaternion multiplication \(PQ\);

\begin{align*}
PQ &= \text{Scalar component}(PQ) + \text{Vector component}(PQ)  \\
   &= \bigr(P_{0}Q_{0} - \vec{P} \cdot \vec{Q}\bigr) \, + \, \bigr(P_0\vec{Q} + Q_0\vec{P} + (\vec{P}\times\vec{Q})\bigr)
\end{align*}


\subsection{What if \(P\) and \(Q\) were pure quaternions?}

If \(P\) and \(Q\) were pure quaternions, then \(P_0\) and \(Q_0\) would both be equal to zero. In that case we would
get;

\begin{align*}
PQ & = \bigl(P_{0}Q_{0} - \vec{P} \cdot \vec{Q}\bigr) \, + \, \bigr(P_0\vec{Q} + Q_0\vec{P} + (\vec{P}\times\vec{Q})\bigr)  \\
   & = \bigl((0 \times 0) - \vec{P} \cdot \vec{Q}\bigr) \, + \, \bigr((0 \times \vec{Q}) + (0 \times \vec{P}) + (\vec{P} \times \vec{Q})\bigr)  \\
   & = \bigl(0 - \vec{P} \cdot \vec{Q}\bigr) \, + \, \bigr(0 + 0 + (\vec{P}\times\vec{Q})\bigr)  \\
   & = -(\vec{P} \cdot \vec{Q}) \, + \, (\vec{P} \times \vec{Q})
\end{align*}

Notice that three of the terms have disappeared and we are now left with just two terms -- the negative of the dot
product between \(P\) and \(Q\), and the cross product between \(P\) and \(Q\).  \\

Let's perform an example multiplication between two pure quaternions, \(\vec{P}\) and \(\vec{Q}\), to see how it's done.
First, let us define \(\vec{P}\) and \(\vec{Q}\) as follows;

\begin{align*}
    P & = \qi 1 + \qj 2 + \qk 3  \\
    Q & = \qi 4 + \qj 5 + \qk 6
\end{align*}

% i(P2 Q3 − P3 Q2 ) + j(P3 Q1 − P1 Q3 ) + k(P1 Q2 − P2 Q1 )

\begin{align*}
    P \times Q = &-\bigl((1 \times 4) + (2 \times 5) + (3 \times 6)\bigr) + \\
                 & \; \qi \bigl((2 \times 6) - (3 \times 5)\bigr) +
                   \qj \bigl((3 \times 4) - (1 \times 6)\bigr) +
                   \qk \bigl((1 \times 5) - (2 \times 4)\bigr)  \\
               = & -\bigl(4 + 10 + 18\bigr) +
                   \qi \bigl(12 - 15\bigr) +
                   \qj \bigl(12 - 6\bigr) +
                   \qk \bigl(5  - 8\bigr)  \\
               = & -32 + \qi (-3) + \qj (6) + \qk (-3) \\
               = & -32 - \qi 3 + \qj 6 - \qk 3
\end{align*}

Note that even though both \(P\) and \(Q\) are 3-dimensional vectors, the result is a 4-dimensional quaternion.